\begin{lemma}[Contained extraction on equal or consecutive parts]
Let $0<\eta<1$, $0<c<d^2/2$, and $d\leq1$, in extended
nonnegative reals. There are $0<\theta<\eta$ and an integer $N\geq1$,
depending only on these parameters, such that the following holds.
If $G$ is a finite simple graph, $U,W$ are disjoint, both sizes belong
to $\{m,m+1\}$ with $m\geq N$, and $(U,W)$ is lower
$(\theta,d)$-super-regular, then there is a graph $F$ on the original
vertex set, contained in $G$, with edges only between $U$ and $W$, for
which $(U,W)$ is $(\eta,c)$-super-regular. Neither coverage of the
ambient vertex set nor absence of other edges in $G$ is required.
\end{lemma}
\begin{proof}
The hypotheses imply that all three parameters are finite,
$0<d\leq1$ and $0<c<1/2$. Work with their real values. Put
$q=\sqrt{2c}$, $g=(d-q)/4$, and $\tau=\eta/4$.
Then $0<q<d\leq1$ and $g>0$. By \noderef{SuperRegularSubgraph},
choose $0<\alpha<\tau$ and a cutoff $N_1$ for extraction with
input density $q$ and output error $\tau$. The same parameters may
be required for two vertex universes at once: take the minimum of
their two supplied errors and the maximum of their two cutoffs.
Weakening an input error only strengthens the required degree and
subset hypotheses, so each supplied extraction still applies.
Choose $0<\theta<\min\{\eta,\alpha/4,g\}$ and an integer
$N\geq\max\{1,N_1,2/\alpha,1/(\alpha g),8/\eta\}$.
These choices precede $G,U,W,m$.

First restrict $G$ to its crossing edges on $U\cup W$, and regard
this graph as a graph on the subtype $U\cup W$. Cardinalities,
opposite-part degrees, and densities of subsets of the two parts are
unchanged. This follows directly from the ordered-edge count in
\noderef{PairDensity}: inclusion of the subtype is a bijection on
each part and on its ordered crossing edges.
The hypotheses of \noderef{LowerSuperRegularPair} give degree
at least $(d-\theta)$ times the opposite size and density at least
$d-\theta$ whenever both subset sizes exceed their $\theta$ thresholds.

If the two sizes are equal, tests above the $\alpha$ thresholds
are above the $\theta$ thresholds, and $d-\theta>q$.
Thus the restricted graph meets the minimum-degree and density
hypotheses of \noderef{SuperRegularSubgraph}: its lower
$(\alpha,q+\alpha)$-regularity is exactly density at least $q$
on those tests, by \noderef{LowerRegularPair}. Extraction yields
error $\tau$ and density $q^2/2=c$. Weakening the error to $\eta$
preserves all degree inequalities and makes the tested family smaller,
by \noderef{RegularPair} and \noderef{SuperRegularPair}.

Suppose next the sizes are $m$ and $h=m+1$, with $U$ the smaller
part. Add one new vertex to $U$ and join it to every vertex of $W$.
The padded parts cover their enlarged vertex set and have size $h$.
Original vertices of $U$ have degree at least $(d-\theta)h>qh$;
the new vertex has degree $h\geq qh$. A vertex of $W$ has degree
at least $(d-\theta)m+1\geq q(m+1)$, since $d-\theta>q$ and $q<1$.
For padded tests $S,T$ of sizes greater than $\alpha h$, put
$S_0=S\setminus\{*\}$. Then
$|S_0|>\alpha(m+1)-1\geq\alpha m/2>\theta m$
and $|T|>\alpha h>\theta h$. The original lower-regularity
test applies, and, even discarding the new edges,
\[
 d_{\rm pad}(S,T)\geq(d-\theta)|S_0|/|S|
 \geq d-\theta-1/|S|>d-\theta-1/(\alpha h)
 \geq d-\theta-g>q.
\]
Here $0<d-\theta\leq1$ and the choice of $N$ were used.
Apply \noderef{SuperRegularSubgraph} to the padded graph.

Delete the new vertex from the extracted graph. Degrees from $U$
to $W$ are unchanged. Degrees from $W$ to $U$ lose at most one,
and therefore lie between $(c-\tau)(m+1)-1$ and
$(c+\tau)(m+1)$. Their error relative to $m$ is at most
$\tau+2/m\leq\eta/2<\eta$, since $c<1/2$, $\tau<1/4$,
and $m\geq8/\eta$. Tests of sizes greater than $\eta m$ and
$\eta(m+1)$ exceed the padded $\tau(m+1)$ thresholds:
$\eta m>\tau(m+1)$ for $m\geq1$. Their edges are unchanged
by deletion, so their densities lie within $\tau$ of $c$ and hence
within $\eta$. These are precisely the clauses of
\noderef{SuperRegularPair}. If $W$ was the smaller part, interchange
the parts in this argument; ordered-edge density is symmetric because
swapping coordinates bijects the two ordered-edge sets.

Finally extend the extracted graph by isolated vertices to the original
ambient type. Its edges are original crossing edges of $G$, and all
tested counts are unchanged. No dummy vertex remains. The real and
extended inequalities agree because every quantity is finite. Input
$d-\theta$ is positive. An output real lower endpoint that is negative
corresponds to truncated endpoint zero and is automatic; otherwise it
is the identical lower bound just proved. This proves the extended
statement as well, including both degree orientations.
\end{proof}
